    \section{Additional Notations}\label{app:notation}
% \begin{align*}
%     &\pi_k ~\text{denotes the permutation in the i-th epoch.} \\
%     &\mA^{(k)}:=(\mI -\eta \vx_{\pi_k(N-1)}\vx_{\pi_k(N-1)}^{\top})\cdots(\mI -\eta \vx_{\pi_k(0)}\vx_{\pi_k(0)}^{\top}) \\
%     &\forall k, \E\mA^{(k)}:=\E\mA = (\mI-\eta\mH)^{N} \\
%     &\mZ^{(k)}_{a \to b} := (\mI -\eta \vx_{\pi_k(a)}\vx_{\pi_k(a)}^{\top})\cdots(\mI -\eta \vx_{\pi_k(b)}\vx_{\pi_k(b)}^{\top}),~~\text{Specifically, $\mA^{(k)} = \mZ^{(k)}_{N-1 \to 0}$}\\
%     &\text{always use }~~K ~~\text{instead of}~~k~~\\ 
%     &\mT^{(k)}:=\mA^{(K)}\cdots\mA^{(k+1)} ~\text{and} ~\mT^{(K)} = \mI \\
%     &\fakeE\mT^{(k)}:=(\mI-\eta\mH)^{K-1-k} \\
%     &\fakeE S_{0 \to k-1}(\mA) := \mI + \E[\mA] + \E[\mA]^2 + \dots + \E[\mA]^{k-1}\\
%     &\mS_{0 \to k-1} := \mI+\mA+\cdots+\mA^{k-1}\\
%     &\mS_l^{(ij)} := \mZ_{N-1 \to \pi_i(l)+1}^{(i)}\vx_l\vx_l^{\top} \mZ_{\pi_j(l)+1 \to N-1}^{(j)}\\
%     &\fakeE \mS_l^{(ij)}:=\fakeE \mS_l^{(ij)} \\
%     &\mS^{(ij)} := \sum_{l=0}^{N-1}\mS_l^{(ij)} \\
%     &\vb^{(k)} := \eta \sum_{l=0}^{N-1}(\mI -\eta \vx_{\pi_k(N-1)}\vx_{\pi_k(N-1)}^{\top})\cdots(\mI -\eta \vx_{\pi_k(l+1)}\vx_{\pi_k(l+1)}^{\top})\xi_{\pi_k(l)}\vx_{\pi_k(l)} \\
%     &= \eta \sum_{l=0}^{N-1}(\mI -\eta \vx_{\pi_k(N-1)}\vx_{\pi_k(N-1)}^{\top})\cdots(\mI -\eta \vx_{\pi_k(\pi_k^{-1}(l)+1)}\vx_{\pi_k(\pi_k^{-1}(l)+1)}^{\top})\xi_{l}\vx_{l} \\
%     &\text{\tingkai{We need to modify the definition of $\mZ^{(k)}_{a \to b}$.}} \\
%     &\mS_{1 \to k-1} := \mA+\cdots+\mA^{k-1} = \mS_k~~~\text{for short}\\
%     &\mS_{i \to j}(\mA):=\mA^{i}+\cdots+\mA^{j}\\
%     & \prod_{l=j}^i \mA^{(l)}:= \mA^{(j)}\mA^{(j-1)}\cdots\mA^{(i)}\;for\;i\le j. \\
%     &\Tilde{\mA}:=\mS\mA\mS^{\top} \text{for a matrix} \mA.
% \end{align*}
% \begin{align*}
%     &\mA^{(k)}:=(\mI -\eta \vx_{\pi_k(N-1)}\vx_{\pi_k(N-1)}^{\top})\cdots(\mI -\eta \vx_{\pi_k(0)}\vx_{\pi_k(0)}^{\top}) \\
%     &\mZ^{(k)}_{a \to b} := (\mI -\eta \vx_{\pi_k(a)}\vx_{\pi_k(a)}^{\top})\cdots(\mI -\eta \vx_{\pi_k(b)}\vx_{\pi_k(b)}^{\top}),~~\text{Specifically, $\mA^{(k)} = \mZ^{(k)}_{N-1 \to 0}$.}\\
%     &\mT^{(k)}:=\mA^{(K)}\cdots\mA^{(k+1)} ~\text{and} ~\mT^{(K)} = \mI. \\
%     &\mS_l^{(ij)} := \mZ_{N-1 \to \pi_i^{-1}(l)+1}^{(i)}\vx_l\vx_l^{\top} \mZ_{\pi_j^{-1}(l)+1 \to N-1}^{(j)}\\
%     &\vb^{(k)} := \eta \sum_{l=0}^{N-1}(\mI -\eta \vx_{\pi_k(N-1)}\vx_{\pi_k(N-1)}^{\top})\cdots(\mI -\eta \vx_{\pi_k(l+1)}\vx_{\pi_k(l+1)}^{\top})\xi_{\pi_k(l)}\vx_{\pi_k(l)} \\
%     &= \eta \sum_{l=0}^{N-1}(\mI -\eta \vx_{\pi_k(N-1)}\vx_{\pi_k(N-1)}^{\top})\cdots(\mI -\eta \vx_{\pi_k(\pi_k^{-1}(l)+1)}\vx_{\pi_k(\pi_k^{-1}(l)+1)}^{\top})\xi_{l}\vx_{l} \\
% \end{align*}
In this section, we provide some additional notations appeared in the following proof of our main results.

\paragraph{Key Quantities.} We define the following key quantities to analyze the sequential updates. For each epoch $k$, let 
\begin{equation}
    \mA^{(k)} := \prod_{i= N-1}^0 (\mI - \eta \vx_{\pi_k(i)}\vx_{\pi_k(i)}^{\top})\label{app:A_def}
\end{equation}
represent the product of sequential updates through all samples in epoch $k$. More generally, we define the partial product operator:
\begin{align*}
    \mZ^{(k)}_{a \to b} &:= \prod_{i=a}^b(\mI - \eta \vx_{\pi_k(i)}\vx_{\pi_k(i)}^{\top}), \quad \text{with} \quad \mA^{(k)} = \mZ^{(k)}_{N-1 \to 0}.
\end{align*}
We further define that $\mZ^{(k)}_{N-1 \to N} = \mI$. The cumulative effect across epochs is captured by:
\begin{align*}
    \mT^{(k)} &:= \prod_{i=K}^{k+1}\mA^{(i)}, \quad \text{and} \quad \mT^{(K)} = \mI.
\end{align*}
% For analyzing the specific contribution of each sample, we define:
% \begin{align*}
%     \mS_l^{(ij)} &:= \mZ_{N-1 \to \pi_i^{-1}(l)+1}^{(i)}\vx_l\vx_l^{\top} \mZ_{\pi_j^{-1}(l)+1 \to N-1}^{(j)},
% \end{align*}
% which isolates the effect of sample $l$ between epochs $i$ and $j$. The accumulation term of noise is given by:
% \begin{align*}
%     \vb^{(k)} &:= \eta \sum_{l=0}^{N-1} (\mI - \eta \vx_{\pi_k(N-1)}\vx_{\pi_k(N-1)}^{\top}) \cdots (\mI - \eta \vx_{\pi_k(l+1)}\vx_{\pi_k(l+1)}^{\top}) \xi_{\pi_k(l)}\vx_{\pi_k(l)} \\
%     &= \eta \sum_{l=0}^{N-1} (\mI - \eta \vx_{\pi_k(N-1)}\vx_{\pi_k(N-1)}^{\top}) \cdots (\mI - \eta \vx_{\pi_k(\pi_k^{-1}(l)+1)}\vx_{\pi_k(\pi_k^{-1}(l)+1)}^{\top}) \xi_{l}\vx_{l},
% \end{align*}
% representing the propagated noise contribution through the update sequence.
% We define a constant $\deltaA:= \Tilde{C}8eD^4\log d$, which is independent of training time $T$ and learning rate $\eta$, and $\Tilde{C}$ is some absolute constant stated in \Cref{lem:un-shuffle-A-EA-expectation}.

\paragraph{Pseudo-expectation Notation $\fakeE$.} Because matrix multiplication is non-commutative and the shuffling in training introduces statistical dependence, the expectations of the random matrices defined above cannot be written in a tractable closed form. To approximate the population excess risk, we therefore introduce the auxiliary notation $\fakeE$.  By construction, $\fakeE$ computes the expectation of each factor as if the variables were independent, deliberately neglecting the correlations.  We then invoke matrix-concentration inequalities to bound the gap between this ``pseudo''-expectation and the true expectation of the original dependent random variables. Specifically, for the above random matrices used in our proof, here we further define that
\begin{align}
  \fakeE\mZ_{a \to b}^{(k)}
    &:= (\mI - \eta \mH)^{a-b+1},\\
  \fakeE\mA^{(k)}
    &:= (\mI - \eta \mH)^{N},\\
  \fakeE\mT^{(k)}
    &:= (\mI - \eta \mH)^{N(K-k)},\label{equ:faek-T}\\
  \fakeE\mS_l^{(ij)}
    &:= \fakeE\!\Bigl[\mZ_{N-1 \to \pi_i^{-1}(l)+1}^{(i)}\Bigr]\,
        \E\!\bigl[\vx_l\vx_l^{\top}\bigr]\,
        \fakeE\!\Bigl[\mZ_{\pi_j^{-1}(l)+1 \to N-1}^{(j)}\Bigr]. \label{equ:faek-S}
\end{align}



% In general, for a random matrix $\mM:=\sum_{i}\prod_{j}^{f(i)}{\mM_{ij}}$, where $\mM_{ij}$ we define 
% \begin{align*}
%     \fakeE\mM:=\sum_{i}\prod_{j}^{f(i)}\E{\mM_{ij}}
% \end{align*}
% to be the fake expectation of $\mB$, which is an approximation of $\E\mB$. 